% =============================================================================
% Rozdział 4: Implementacja
% =============================================================================
\chapter{Implementacja}

Rozdział niniejszy dokumentuje implementację systemu badawczego. Wszystkie listingi odwołują się do rzeczywistych plików w repozytorium projektu.

\section{Struktura repozytorium}

Kod projektu zorganizowano w następującej strukturze:

\begin{verbatim}
magisterka/
  apps/
    cpu-service/         # CPU-bound microservice
    io-service/          # I/O-bound microservice
    echo-server/         # Auxiliary echo service
    memory-service/      # Memory-bound microservice (extension)
  k8s/
    base/                # Namespace
    cpu-service/         # Deployment + Service + HPA x3
    io-service/          # Deployment + Service + HPA x3
    echo-service/        # Deployment + Service
    memory-service/      # Deployment + Service + HPA
    monitoring/          # Prometheus Adapter config
  monitoring/
    dashboards/          # Grafana dashboard JSON
  scripts/
    run-tests.sh         # Test orchestration (8 scenarios)
  tests/
    k6/                  # k6 load test scripts
  docs/
    prd.md               # Product Requirements Document
  docker-compose.yaml    # Optional local development (not used in research)
\end{verbatim}

\section{Implementacja mikroserwisów}

Każdy mikroserwis jest niezależną aplikacją FastAPI (Python 3.12), konteneryzowaną przez Docker i uruchamianą przez Uvicorn. Wszystkie serwisy współdzielą wspólny wzorzec implementacyjny: obraz bazowy \texttt{python:3.12-slim}, instalację zależności przez \texttt{pip install --no-cache-dir}, automatyczną instrumentację Prometheus przez \texttt{prometheus-fastapi-instrumentator} (metryki HTTP na endpointcie \texttt{/metrics}) oraz endpoint \texttt{GET /health} wykorzystywany przez sondy Kubernetes.

\begin{lstlisting}[language=Dockerfile, caption={Dockerfile --- wzorzec wspólny dla wszystkich serwisów}, label={lst:dockerfile}]
FROM python:3.12-slim
WORKDIR /app
COPY requirements.txt .
RUN pip install --no-cache-dir -r requirements.txt
COPY main.py .
EXPOSE 8000
CMD ["uvicorn", "main:app", "--host", "0.0.0.0", "--port", "8000"]
\end{lstlisting}

Obrazy są budowane lokalnie i przesyłane do GCP Artifact Registry, skąd GKE pobiera je automatycznie podczas wdrażania na podstawie pola \texttt{image} w Deployment.

\subsection{CPU-service}

CPU-service implementuje dwa endpointy obciążające procesor. \texttt{POST /process} akceptuje obraz w formacie multipart form-data i wykonuje potok przetwarzania CPU-intensywnego z biblioteką Pillow: konwersja do skali szarości, filtry BLUR, SHARPEN i FIND\_EDGES, ekstrakcja kanałów RGB z resize. Po przetworzeniu, dla każdego kanału obliczana jest entropia Shannona:

\begin{equation}
  H = -\sum_{i=0}^{255} p_i \log_2(p_i), \quad \text{gdzie } p_i = \frac{\text{histogram}[i]}{\text{total\_pixels}}
\end{equation}

\texttt{GET /fibonacci} oblicza $n$-tą liczbę Fibonacciego rekurencyjnym algorytmem o złożoności $O(2^{n})$, co gwarantuje wysokie zużycie CPU przy niewielkiej liczbie żądań.

\lstinputlisting[
  language=Python,
  caption={Rekurencyjna implementacja ciągu Fibonacciego},
  label={lst:cpu-fib},
  firstline=82,
  lastline=90
]{../apps/cpu-service/main.py}

Serwis wystawia również niestandardowe metryki Prometheus: \texttt{cpu\_service\_images\_processed\_total} (Counter) oraz \texttt{cpu\_service\_processing\_duration\_seconds} (Histogram).

\subsection{IO-service}

IO-service symuluje operacje I/O przez dwa endpointy. \texttt{GET /query} wykonuje sekwencję wywołań \texttt{asyncio.sleep()} z konfigurowalnym opóźnieniem (\texttt{delay}, domyślnie 100\,ms), liczbą kroków (\texttt{steps}, domyślnie 3) i losowym rozrzutem (\texttt{jitter}, $\pm20\%$). Całkowity czas odpowiedzi wynosi $\text{steps} \times \text{delay} \times (1 \pm \text{jitter})$, a procesor pozostaje bezczynny podczas oczekiwania.

\lstinputlisting[
  language=Python,
  caption={Endpoint GET /query w io-service},
  label={lst:io-query},
  firstline=195,
  lastline=252
]{../apps/io-service/main.py}

\texttt{GET /upstream} przekazuje żądanie do echo-server przez \texttt{httpx.AsyncClient}. Serwis wykorzystuje \texttt{asyncio.Semaphore(50)} do ograniczenia współbieżności oraz wystawia metryki: \texttt{io\_service\_queries\_total} (Counter), \texttt{io\_service\_upstream\_duration\_seconds} (Histogram) i \texttt{io\_service\_concurrent\_queries} (Gauge).

\subsection{Echo-server i Memory-service}

Echo-server jest minimalistycznym serwisem pomocniczym (endpointy \texttt{GET /health} i \texttt{POST /echo}), wyłączonym z macierzy badawczej. Memory-service implementuje pamięciową bazę klucz-wartość z limitem 10000 wpisów i endpointem badawczym \texttt{POST /cache/fill} do wypełniania cache danymi testowymi o kontrolowanym rozmiarze. Nie został włączony do głównej macierzy badawczej.

\section{Konfiguracja Kubernetes}

\subsection{Deploymenty}

Każdy Deployment definiuje: początkową liczbę replik (2), politykę pobierania obrazów (\texttt{imagePullPolicy: Always}), sondy liveness i readiness (\texttt{GET /health} co 10\,s, initialDelay 5\,s) oraz resource requests i limits zgodnie z Tabelą~\ref{tab:resource-allocation} (Rozdział~3).

\lstinputlisting[
  language=yaml,
  caption={Deployment cpu-service},
  label={lst:deployment-cpu},
  firstline=1,
  lastline=45
]{../k8s/cpu-service/deployment.yaml}

\subsection{HPA --- trzy warianty}

Dla cpu-service i io-service zdefiniowano po trzy warianty HPA (API \texttt{autoscaling/v2}, zakres 2--10 replik), różniące się wyłącznie typem metryki. HPA CPU używa progu 50\% średniego zużycia CPU. HPA Memory używa progu 70\% średniego zużycia pamięci. HPA Custom RPS używa progu 100 żądań na sekundę na Pod, z metryką dostarczaną przez prometheus-adapter.

\lstinputlisting[
  language=yaml,
  caption={HPA CPU dla cpu-service},
  label={lst:hpa-cpu},
  firstline=1,
  lastline=22
]{../k8s/cpu-service/hpa-cpu.yaml}

\lstinputlisting[
  language=yaml,
  caption={HPA Memory dla cpu-service},
  label={lst:hpa-memory},
  firstline=1,
  lastline=23
]{../k8s/cpu-service/hpa-memory.yaml}

\lstinputlisting[
  language=yaml,
  caption={HPA Custom RPS dla cpu-service},
  label={lst:hpa-custom},
  firstline=1,
  lastline=24
]{../k8s/cpu-service/hpa-custom.yaml}

\subsection{Services}

Wszystkie serwisy aplikacyjne są typu \texttt{LoadBalancer} --- GKE automatycznie tworzy zewnętrzny GCP LoadBalancer (Layer 4, TCP) i przydziela mu publiczny adres IP, na który kierowany jest ruch z maszyny k6-generator.

\section{Infrastruktura monitorująca}

\subsection{kube-prometheus-stack}

Pakiet kube-prometheus-stack jest instalowany przez Helm, co zapewnia automatyczną konfigurację i spójność wersji wszystkich komponentów (Prometheus, Grafana, kube-state-metrics, node-exporter, Prometheus Operator). Parametr \texttt{grafana.service.type=LoadBalancer} eksponuje Grafanę publicznie, umożliwiając obserwację eksperymentów na żywo.

\subsection{Prometheus Adapter}

Prometheus Adapter v0.12.0 wdrażany jest przez osobne manifesty YAML (nie jest częścią pakietu Helm). Jego konfiguracja w ConfigMap definiuje reguły PromQL mapujące metryki Prometheusa na format Custom Metrics API:

\lstinputlisting[
  language=yaml,
  caption={Konfiguracja Prometheus Adapter},
  label={lst:adapter-config-impl},
  firstline=1,
  lastline=66
]{../k8s/monitoring/prometheus-adapter-config.yaml}

\subsection{Grafana --- dashboard i eksport}

Autorski dashboard ,,Magisterka -- System Overview'' zawiera 16 paneli odświeżanych co 10 sekund, pokazujących RPS, zużycie CPU i pamięci per Pod, percentyle opóźnień, liczbę replik i stan HPA. Dla każdego scenariusza eksportowanych jest 5 zestawów CSV: \texttt{replicas.csv}, \texttt{cpu\_mcpu.csv}, \texttt{memory\_mb.csv}, \texttt{rps.csv} i \texttt{latency\_ms.csv}.

\section{Skrypty automatyzujące}

\subsection{run-tests.sh --- orkiestracja eksperymentów}

Skrypt \texttt{scripts/run-tests.sh} automatyzuje wykonanie wszystkich 8 scenariuszy badawczych. Działa na maszynie k6-generator (Compute Engine), która ma skonfigurowany dostęp do klastra GKE przez \texttt{kubectl}. Dla każdego scenariusza skrypt: przełącza HPA na odpowiednią strategię (funkcja \texttt{switch\_hpa}), resetuje Deployment do 1 repliki z oczekiwaniem na gotowość, uruchamia test k6 z ruchem kierowanym na adresy IP GCP LoadBalancer, zapisuje wyniki do \texttt{results/<scenario>/}. Po każdym scenariuszu następuje cooldown: skalowanie obu serwisów do 1 repliki, usunięcie HPA, odczekanie 120 sekund.

\lstinputlisting[
  language=bash,
  caption={Funkcja switch\_hpa() --- przełączanie strategii HPA},
  label={lst:switch-hpa},
  firstline=137,
  lastline=184
]{../scripts/run-tests.sh}

\lstinputlisting[
  language=bash,
  caption={Funkcja cooldown() --- stabilizacja klastra między scenariuszami},
  label={lst:cooldown},
  firstline=202,
  lastline=220
]{../scripts/run-tests.sh}

\section{Scenariusze testowe k6}

Każdy skrypt k6 definiuje scenariusze obciążeniowe, progi akceptacji oraz mapowanie hostów na adresy IP GCP LoadBalancer.

\lstinputlisting[
  language=javascript,
  caption={Scenariusze i progi w cpu-load-test.js},
  label={lst:k6-cpu-scenarios},
  firstline=41,
  lastline=72
]{../tests/k6/cpu-load-test.js}

\lstinputlisting[
  language=javascript,
  caption={Scenariusze i progi w io-load-test.js},
  label={lst:k6-io-scenarios},
  firstline=36,
  lastline=76
]{../tests/k6/io-load-test.js}

\section{Podsumowanie implementacji}

Zaimplementowany system badawczy składa się z czterech mikroserwisów, infrastruktury monitorującej opartej na kube-prometheus-stack (Helm) oraz zestawu skryptów automatyzujących testy. Kluczowe cechy implementacji: pełna automatyzacja procesu badawczego przez skrypt \texttt{run-tests.sh}, deklaratywna konfiguracja wszystkich zasobów Kubernetes (możliwa do odtworzenia na dowolnym klastrze GKE), monitoring przez Prometheusa i Grafanę z eksportem 5 zestawów CSV na scenariusz oraz wykorzystanie GCP LoadBalancer zamiast port-forward dla realistycznych warunków sieciowych.
